\scanpage{13}
\deleted{2. Let $X$ be a space, $Y$ a closed subspace.}
\unsure{\deleted{(6)}|(6)} \deleted{Let $T$}

\uline{NB} The $T_\lambda$ $(\lambda\in[0,1])$ form
an increasing family of closed subsets of
$T$, the $\mathring T_\lambda\overset{\mathrm{def}}=
\lambda\mathring T\deleted{\unsure{_\lambda}}$ $(\lambda\in]0,1])$
\margin{$\Vert$}
an increasing family of open subsets. If $\lambda\in]0,1[$, then
$\mathring T_\lambda$ is indeed the \inserted{topological} interior
of $T_\lambda$ in $T$,
but of course $\mathring T_1=\mathring T$ is not the
\inserted{\unsure{topological|topological}} interior of $T_1$ in
the topological space $T_1=T$, i.e. we do not have
$\mathring T\unsure{\deleted{_1}|_1|\deleted{_\lambda}}=T$ unless
$\dot T=\phi$ i.e. $Y=T$.

Put, for $\lambda\in]0,1[$
\[
(\unsure{6})\qquad
\left\{\begin{aligned}
T\setminus\mathring T_\lambda
&\overset{\mathrm{def}}=\deleted{\unsure{E|\xi}}
\mathrm{Ep}^{(T)}_{1-\lambda=\varepsilon}
\quad\text{``thickening of index $\varepsilon$''}\\
&\hspace{4em}\text{of the boundary $\dot T$ in $T$}\\
\dot{\mathrm{Ep}}_\varepsilon^{(T)}
&\overset{\mathrm{déf}}=\dot T\unsure{\amalg|\cup}(1-\varepsilon).\dot T
\end{aligned}\right.
\]
\uline{NB} The boundary of $E_\varepsilon$ in $T$ is
$(1-\varepsilon)\dot T$, and $E_\varepsilon$
is a neighborhood of $\dot T$. We have
\[
(7)\qquad\mathrm{Ep}_\varepsilon^{(T)}
\simeq[\underbrace{1-\varepsilon}_{\lambda},1]\times\dot T
\]
\margin{\uline{NB} If $Y,\inserted{\dot T}$ compact, \deleted{\unsure{the}}
then $T$ compact, ? and the $T_\varepsilon$
$(\varepsilon\in]0,1])$ form
a fundamental system of neighborhoods of $Y$ in $T$,
\deleted{\unsure{closed}}\inserted{\unsure{of open sets|open}}.
% crops/150_scan013_P07.png: the following live fragment lies below syst. fond.; its attachment is uncertain.
\unsure{more|full} : $T$
(as tubular \unsure{and} absolute neighborhoods of $Y$)}
% P07 scope check: the nearby overworked (6) and underlined main-text NB belong to P05/P06, not to this marginal note.

\textcircled{2} Let $Y\lhook\joinrel\relbar X$ be a closed inclusion of topological
 spaces. A \uline{conical neighborhood} of \uline{$Y$ in} $X$
is a neighborhood $T$ of $Y$ in $X$
equipped with an action (1) $I\times T\mathrel{\relbar\joinrel\relbar}T$
\inserted{making $T$ an absolute \deleted{tubular}\inserted{conical}
neighborhood of $Y$.} \deleted{satisfying the conditions \unsure{stated.|id.}}
Then the
$\lambda T=T_\lambda$ $(\lambda\in]0,1])$
form other tubular
neighborhoods of $Y$ in $X$ (increasing family) and
we see that for $\lambda\in]0,1[$, we have
\[
(\unsure{8})\qquad\dot T_\lambda=\lambda\dot T
\quad\text{is the boundary of }T_\lambda\unsure{=|:}\lambda T\text{ in }X.
\]
\margin{12}
